\section{The rank 0 case}\label{sec:rank0}

It might be appropriate here to discuss the rank 0 case. The metaplectic
covers of $A(\mathbb{A})$ are not in general abelian but nonetheless, the
trace formula takes an especially simple form in this case. In this section,
we shall prove the trace formula for metaplectic covers of $A({\mathbb{A}})$ and
apply it to prove a ``matching'' between automorphic forms on $A({\mathbb{A}})$
and on metaplectic covers of $A({\mathbb{A}})$. At the same time, we construct
invariant distributions $I_{\mathcal{O}}^A$ and $I_\chi ^A$ which will be
needed in an inductive definition of certain invariant distributions
on $\overline{G({\mathbb{A}})}$ in Section~\ref{sec:invariant}. 

\subsection{Metaplectic lifting over the split torus}\label{ssec:rank0lift}

Suppose that $k/{\mathbb{Q}}$ is a number field containing the $n-th$ roots of
unity. (If $n>2$ then this implies $k$ has no real places, i.e, $k_v=\mathbb{C}$
for all $v\in \Sigma _k^\infty $.)
Let $A$ and $A({\mathbb{A}})^1$ be as in the previous section,
so $A({\mathbb{Q}})=k^\times$.
\emph{We shall often identify }
$A({\mathbb{A}})$
\emph{\ with }
${\mathbb{A}}_k^{\times }$
\emph{\ in this section}. Let
$(..,...)_n:A({\mathbb{A}})\times A({\mathbb{A}})
\cong {\mathbb{A}}_k^{\times }\times {\mathbb{A}}_k^{\times}\rightarrow \mu _n$
denote the global Hilbert symbol defined by
\cite[section XIII.5, proposition 7]{W1} as a product of local
Hilbert symbols $(...,...)_{n,v}$, where
$\mu _n$ is the group of $n$-th roots of unity. Let $\mu _n$ act on 
$G({\mathbb{A}})={\mathbb{A}}_k^{\times }$ trivially.
A 2-cocycle on $A({\mathbb{A}})$ with
values in $\mu _n$ is a continuous mapping
$\beta :A({\mathbb{A}})\times A({\mathbb{A}})\rightarrow \mu _n$ satisfying 
\begin{equation}\label{eq:3.1}
\beta (g_2,g_3)\beta (g_1,g_2g_3)=\beta (g_1,g_2)\beta (g_1g_2,g_3),\;
\end{equation}
$g_1,g_2,g_3\in {\mathbb{A}}_k^{\times }$. We also denote by $\beta $ the
restriction
$\beta :A({\mathbb{A}})^1\times A({\mathbb{A}})^1\rightarrow \mu _n$. This
is a product of local 2-cocycles, denoted $\beta _v$. Let
$\overline{A(\mathbb{A})}^1$ denote the $n$-fold
metaplectic cover of $A({\mathbb{A}})^1$
defined by the 2-cocycle $\beta $ given by 
\begin{equation}\label{eq:3.2}
\beta (g_1,g_2)=(g_1^{-1},g_2^{-1})_n\cdot (-g_1g_2^{-1},g_1^{-1}g_2^{-1})_n.
\end{equation}
The fact that this satisfies the condition above may be verified by
multiplying each term out and using the properties of the Hilbert symbol
listed in \cite[Proposition 60, p. 176]{Sha} (this is long but
straightforward and in any case well-known).
By definition of group extension, multiplication
on $\overline{A({\mathbb{A}})^1}=A({\mathbb{A}})^1\times \mu _n$
(as \emph{sets}) is given by 
\begin{equation}\label{eq:3.3}
(g,\varsigma )\cdot (g^{\prime },\varsigma ^{\prime })=(gg^{\prime },\beta
(g_1,g_2)\varsigma \varsigma ^{\prime }),
\;\;\;\;g,g^{\prime }\in A({\mathbb{A}}),
\;\;\varsigma ,\varsigma ^{\prime }\in \mu _n.
\end{equation}
This is a central extension of $A({\mathbb{A}})^1$ (i.e.,
$\mu _n$ belongs to the
center of the metaplectic cover) which fits into the following exact
sequence 
\[
1\rightarrow \mu _n\rightarrow \overline{A({\mathbb{A}})^1}
\rightarrow A({\mathbb{A}})^1\rightarrow 1,
\]
where $\mu _n\rightarrow \overline{A({\mathbb{A}})^1}$ is the
inclusion of $\mu _n$ into the second component of
$\overline{A({\mathbb{A}})}^1$, and the map
$\overline{A({\mathbb{A}})^1}\rightarrow A({\mathbb{A}})^1$ is the
projection map.

Note that
$\overline{A({\mathbb{A}})}^1=\overline{A({\mathbb{A}})^1}$.

Using the local 2-cocycles $\beta _v$, we may in turn define local n-fold
covers $\overline{k_v^{\times }}$ of $k_v^{\times },\;v\in \Sigma _k$. The
rule for multiplication on the local covers
\[
\overline{A({\mathbb{Q}}_p)}=
\overline{\prod_{v|p}k_v^{\times }}
\subset \prod_{v|p}
\overline{k_v^{\times }}
\]
of $\prod_{v|p}k_v^{\times }$ is the same, except one must
replace $\beta $ by $\prod_{v|p}\beta _v$. The topology on
$\overline{k_v^{\times }}$
is defined by stipulating 
\[
\{U_v\times 1\;|\;\beta _v(U_v,U_v)=1\;\mathrm{and}\;1\in U_v\subset
k_v^{\times }\;\mathrm{is\;open}\}
\]
to be a basis of open neighborhoods of the identity.
If $(p,n)=1$ and if $v|p$ then the n-fold covers
$\overline{A(\mathbb{Z}_p)}$ and
$\overline{{\mathcal{O}}_v^{\times }}$ split by
the properties of the Hilbert symbol listed in \cite[section XIII.5]{W1}.

The global cover $\overline{A({\mathbb{A}})}$ is not the restricted direct
product of the local covers $\overline{k_v^{\times }}$. However, there is an
embedding 
\[
\overline{A({\mathbb{A}})}\hookrightarrow
\prod_{v\in \Sigma _k}^{\prime }\overline{k_v^{\times }}.
\]

We call a representation $\pi $ of $\overline{A({\mathbb{A}})^1}$
\textbf{genuine} if it satisfies

\begin{equation}\label{eq:3.4}
\pi (g,\varsigma )=\varsigma^{-1} \cdot \pi (g,1).
\end{equation}
A {\em function} satisfying this property will sometimes be
called {\bf anti-genuine}.
Similarly, we define a representation $\pi $ of
$\overline{A({\mathbb{A}})}=\overline{{\mathbb{A}}_k^{\times }}$ or
of $\overline{k_v^{\times }}$ to be 
\textbf{genuine} if it satisfies \eqref{eq:3.4}.

If $H<\overline{A({\mathbb{A}})}$ a locally compact totally disconnected group
of type 1, let $\Pi (H)$ denote
the set of equivalence classes of admissible genuine
representations of $H$ and
let $\Pi_u(H)$ denote the subset of unitary ones.
Let $\Pi_f (\overline{A({\mathbb{A}})})$
denote the subset (equivalance classes of) representations in
$\Pi(\overline{A({\mathbb{A}})})$ which can be written as the
restriction to $\overline{A({\mathbb{A}})}$
of a {\em factorizable} representation $\overline{\pi }$
of the restricted direct product 
\[
\Pi (\prod_v\overline{k_v^{\times }})=\prod_v
\Pi (\overline{k_v^{\times }}).
\]
Let $\Pi_f (\overline{A({\mathbb{A}})}^1)$ denote the
set of equivalence classes of
representations of $\overline{A({\mathbb{A}})^1}$ which can be written as the
restriction to $\overline{A({\mathbb{A}})^1}$ of an element of
$\Pi_f (\overline{A({\mathbb{A}})})$.
Let $\Pi_{u,f} (\overline{A({\mathbb{A}})})$,
$\Pi_{u,f} (\overline{A({\mathbb{A}})^1})$
denote the corresponding subset of factorizable unitary representations.

\begin{lemma}\label{lem:3.1} (\cite[\S 3]{Ge2})
\[
\begin{array}{c}    
\Pi_f (\overline{A({\mathbb{A}})})=\Pi (\overline{A({\mathbb{A}})}),\\
\Pi_{u,f} (\overline{A({\mathbb{A}})})=\Pi_u (\overline{A({\mathbb{A}})}),\\
\Pi_f (\overline{A({\mathbb{A}})^1})=\Pi (\overline{A({\mathbb{A}})^1}),\\
\Pi_{u,f} (\overline{A({\mathbb{A}})^1})=\Pi_u (\overline{A({\mathbb{A}})^1})
\end{array}
\]
\end{lemma}

See also \cite{Ge1}, \cite{F} for the non-metaplectic version.

We shall sketch a proof of the following

\begin{theorem}\label{thm:3.2}
There is a natural 1-1 correspondence between the subset of all genuine
irreducible representations
$\overline{\chi }\in \Pi (\overline{A({\mathbb{A}})}^1)$
of $\overline{A({\mathbb{A}})^1}$ and characters of the form
$\chi ^{n/2}$ (n even), $\chi ^n$ (n odd),
for $\chi \in $\ $\Pi (A({\mathbb{A}})^1)$.
\end{theorem}

First we prove a trace formula for $\overline{A({\mathbb{A}})}^1$.

\subsection{Irreducible representations of $\overline{A({\mathbb{A}})}^1$}

For any group $H$, let
Let 
\begin{equation}\label{eq:3.5}
H^m
=\{x\in H\;|\;x=y^m,\;\;\mathrm{some}
\;\;y\in H\}.
\end{equation}
Let 
\begin{equation}\label{eq:3.6}
N=
\left\{
\begin{array}{c}
n,\ \ n \mathrm{\ odd},\\
n/2,\ \ n \mathrm{\ even}.
\end{array}
\right.
\end{equation}

We have, by the properties of the Hilbert
symbol,
$\overline{A({{\mathbb{A}}})^{1,n}}=A({{\mathbb{A}}})^{1,n}\times \mu _n$,
that is to say, the metaplectic cover
splits over $A({{\mathbb{A}}})^{1,n}$.
By Lemma~\ref{lem:2.6}(e) above, the cover
$\overline{A({{\mathbb{A}}})^{1,N}}$ is abelian (though it
does not split).

For $\chi \in \Pi (\overline{A({\mathbb{A}})^1})$,
$h\in \overline{A({\mathbb{A}})}$, define
$\chi ^h(g)=\chi (h^{-1}gh)$, for $g\in 
\overline{A({\mathbb{A}})}^1$. We call such a
$\chi $ \textbf{self-conjugate} if
there is an $h$ as above such that $\chi =\chi ^h$
in $\Pi (\overline{A({\mathbb{A}})^1})$.
If $h=(x,\varsigma ),\;\varsigma \in \mu _n,$ then 
\[
\chi ^h(g)=\chi (h^{-1}gh)
=\chi (x^{\prime },(x,x^{\prime })_n^2\varsigma^{\prime }),
\;\;\;\;\;\;\;\;\;\;g=(x^{\prime },
\varsigma ^{\prime })\in \overline{A({\mathbb{A}})^1},
\]

If $p<\infty$, let $C_p\subset A({\mathbb{Q}}_p)$ denote a subgroup
of finite index such that
$\overline{C_p}\subset \overline{A({\mathbb{Q}}_p)}$ is
an abelian subgroup. (In particular, we may take
$C_p=A({\mathbb{R}})^N$.) If $p=\infty$ let $C_p=A({\mathbb{R}})$. Let
\begin{equation}\label{eq:3.7}
C_{{\mathbb{A}}}=\prod_{p\leq \infty }C_p
\end{equation}
denote the restricted direct product.

We have the following result (due to Clifford, Mackey,
...)

\begin{lemma}\label{lem:3.3}
If $\chi \in \Pi (\overline{C_{{\mathbb{A}}}})$ is not self-conjugate then 
\[
Ind_{\overline{C_{{\mathbb{A}}}}}^{\overline{A({\mathbb{A}})}}(\chi )
\cong Ind_{\overline{C_{{\mathbb{A}}}}}^{\overline{A({\mathbb{A}})}}(\chi ^h),
\;\;\;h\in  \overline{A({\mathbb{A}})},
\]
is irreducible and belongs to $\Pi (\overline{A({\mathbb{A}})})$.
\end{lemma}

\begin{remark}\label{rem:3.4} By analogy with Mackey's irreduciblity
criterion in \cite[\S 7.4]{Se}, the converse to the avove result also
holds. The detailed proof of the local result can be found in
\cite{J3}.
\end{remark}

\textsf{proof}: We must first reduce this down to a statement about local
representations and then prove the local statement.

Pick a $p\leq \infty $. By Clifford-Mackey theory (\cite{J3}, \cite[\S 7.4]{Se}),
if $\chi \in \Pi (\overline{A({\mathbb{Q}}_p)})$ is not self-conjugate then 
\begin{equation}\label{eq:3.8}
Ind_{\overline{C_p}}^{\overline{A({\mathbb{Q}}_p)}}(\chi )
\cong Ind_{\overline{C_p}}^{\overline{A({\mathbb{Q}}_p)}}(\chi ^h),
\;\;\;h\in \overline{A({\mathbb{Q}}_p)},
\end{equation}
is irreducible and belongs to $\Pi (\overline{A({\mathbb{Q}}_p)})$. We may take
the restricted direct product of these local representations to obtain an
irreducible (by Schur's lemma) factorizable representation on
the restricted direct product of of these local covers. $\Box $

In our case, it is easy to see that for any
$\chi \in \Pi (\overline{C_p})$,
the restriction of
$Ind_{\overline{C_p}}^{\overline{A({\mathbb{Q}}_p)}}(\chi )$
to the abelian subgroup
$\overline{C_p}$ must decompose into a direct sum of
characters: 
\begin{equation}\label{eq:3.9}
Ind_{\overline{C_p}}^{\overline{A({\mathbb{Q}}_p)}}(\chi )|_{\overline{A_p}}
\cong \bigoplus_{i=1}^dm_i\chi _i,
\;\;\;\chi _i\in \Pi (\overline{C_p}),
\end{equation}
where $m_i,d$ is are some integers and
one of these $\chi _i$ must be $\chi $ itself. By the isomorphism
(for any $h\in \overline{A({\mathbb{Q}}_p)}$)
\begin{equation}\label{eq:3.10}
Ind_{\overline{C_p}}^{\overline{A({\mathbb{Q}}_p)}}(\chi )
\cong Ind_{\overline{C_p}}^{\overline{A({\mathbb{Q}}_p)}}(\chi ^h),
\;\;\;h\in \overline{A({\mathbb{Q}}_p)},
\end{equation}
one of these $\chi _i$ must be equal to $\chi ^h$.
In fact, we have

\begin{lemma}\label{lem:3.5} (\cite[Lemma 2.1]{GeKn})
If $\chi \in \Pi (\overline{C_p})$ then all the $m_i's$ in \eqref{eq:3.9} are
equal (to $m$, say) and
\[
Ind_{\overline{C_p}}^{\overline{A({\mathbb{Q}}_p)}}(\chi )|_{\overline{C_p}}
\cong \bigoplus_{h\in \overline{A({\mathbb{Q}}_p)}
/\overline{A({\mathbb{Q}}_p)}_{\chi}}
m\chi ^h,
\]
where
\[
\overline{A({\mathbb{Q}}_p)}_{\chi}
=\{a\in\overline{A({\mathbb{Q}}_p)} \ |\ \chi^a=\chi\}.
\]
\end{lemma}

\begin{remark}\label{rem:3.6} By analogy with \cite[ch 7]{Se}, one might expect there to
be a more general form of this result, closely related to
Frobenius reciprocity.
\end{remark}
\begin{remark}\label{rem:3.7}
An aside: Even using the explicit local recipocity law of
(Sen, Wiles, ...) it does not appear to be
easy to explicitly determine $C_p$
such that $A({\mathbb{Q}}_p)/C_p$ is minimal.
\end{remark}


\subsection{The trace formula for \emph{$\overline{A({\mathbb{A}})^1}$}}

There is natural projection
$pr:\overline{A({\mathbb{A}})}\;\rightarrow \;A({\mathbb{A}}),
\;(x,\varsigma )\longmapsto x$. Let
${\mathcal{S}}(\overline{A({\mathbb{A}})}^1)$
denote the collection of all genuine functions $f$ on
$\overline{A({\mathbb{A}})^1}$ such
that $f\circ pr\in {\mathcal{S}}(A({\mathbb{A}})^1)$. Let 
\[
H=
L^2(\overline{A({\mathbb{Q}})}\backslash \overline{A({\mathbb{A}})^1)})_{gen},
\]
where \emph{$\overline{A({\mathbb{Q}})}$} is regarded as a subgroup
of $\overline{A({\mathbb{A}})^1}$ via the
diagonal embedding. Let $R$ denote
the right regular representation: 
\[
R:\overline{A({\mathbb{A}})^1}\rightarrow \;End(H),
\]
defined by $(R(y)\varphi )(x)
=\varphi (xy),\;x,y\in \overline{A({\mathbb{A}})}^1,
\;\varphi \in H$.
For $\overline{f}\in {\mathcal{S}}(\overline{A({\mathbb{A}})}^1)$, let 
\[
R(\overline{f})=\int_{\overline{A({\mathbb{A}})}^1}\overline{f}(x)R(x)dx.
\]
Since $R(x):H\rightarrow H$ is a unitary operator (in particular bounded),
this integral converges. This
operator $R(\overline{f}):H\rightarrow H$ may be written 
\begin{equation}\label{eq:3.11}
\begin{array}{c}
(R(\overline{f})\varphi )(x)
=\int_{\overline{A({\mathbb{A}})}^1}\overline{f}(x)(R(y)\varphi )(x)dy,
\;\;\;\varphi \in H \\ 
=\int_{\overline{A({\mathbb{A}})}^1}\overline{f}(x^{-1}y)\varphi (y)dy \\ 
=\int_{A({\mathbb{Q}})\backslash \overline{A({\mathbb{A}})}^1}
\sum_{\delta \in A({\mathbb{Q}})}
\overline{f}(x^{-1}\delta y)\varphi (\delta y)dy \\ 
=\int_{A({\mathbb{Q}})\backslash \overline{A({\mathbb{A}})^1}}
\varphi (y)
\sum_{\delta \in A({\mathbb{Q}})}\overline{f}(x^{-1}\delta y)dy.
\end{array}
\end{equation}
Let 
\[
K(x,y)=K_{\overline{f}}(x,y)
=\sum_{\delta \in A({\mathbb{Q}})}
\overline{f}(x^{-1}\delta y).
\]

For each $\overline{\chi }\in \Pi (\overline{A({\mathbb{A}})}^1)$, let 
\begin{equation}\label{eq:3.12}
\begin{array}{c}
\overline{\chi }(\overline{f})
=\int_{\overline{A({\mathbb{A}})}^1}
\overline{f}(y)\overline{\chi }(y)dy, \\ 
tr\;\overline{\chi }(\overline{f})
=\int_{\overline{A({\mathbb{A}})^1}}\overline{f}(y)
tr\;(\overline{\chi }(y))dy
=\overline{f}^{\wedge }(\overline{\chi }).
\end{array}
\end{equation}
We call the latter integral the \textbf{Fourier transform}
of $\overline{f}$ at $\overline{\chi }$, and the
former integral the \textbf{operator-valued}
\textbf{Fourier transform} of $\overline{f}$ at
$\overline{\chi }$. Note that these are
$0$ if $\overline{f}$ is not genuine.

Finally, for each $h\in $ $\overline{A({\mathbb{A}})^1}$, let
$\overline{A({\mathbb{A}})}_h^1$ denote the \textbf{centralizer} of $h$
in $\overline{A({\mathbb{A}})^1}$: 
\begin{equation}\label{eq:3.13}
\overline{A({\mathbb{A}})}_h^1
=\{A\in \overline{A({\mathbb{A}})}^1\;|\;ghg^{-1}=h\}.
\end{equation}

\begin{theorem}\label{thm:3.8}
(a) $tr\;R(\overline{f})
=\int_{A({\mathbb{Q}})\backslash \overline{A({\mathbb{A}})}^1}
K_{\overline{f}}(x,x)dx.$

(b) For each $\overline{\chi }\in \Pi (\overline{A({\mathbb{A}})}^1)$, there is
an integer $m(\overline{\chi })\geq 0$\ such that 
\[
R(\overline{f})
=\bigoplus_{\overline{\chi }\in \Pi_u (\overline{A({\mathbb{A}})}^1)}
m(\overline{\chi })\overline{\chi }(\overline{f}).
\]
\end{theorem}

Thus we have 
\begin{equation}\label{eq:3.14}
\begin{array}{c}
tr\;R(\overline{f})
=\sum_{\overline{\chi }\in \Pi (\overline{A({\mathbb{A}})}^1)}
m(\overline{\chi })tr(\overline{\chi }(\overline{f})) \\ 
=\int_{A({\mathbb{Q}})\backslash \overline{A({\mathbb{A}})}^1}
\sum_{\delta\in A({\mathbb{Q}})}
\overline{f}(x^{-1}\delta x)dx \\ 
=\int_{A({\mathbb{Q}})\backslash \overline{A({\mathbb{A}})}^1}
\sum_{\{\delta \}}
\sum_{\gamma \in A({\mathbb{Q}})_\delta\backslash A({\mathbb{Q}})}
\overline{f}(x^{-1}\gamma ^{-1}\delta \gamma x)dx \\ 
=\sum_{\{\delta \}}
\int_{A({\mathbb{Q}})\backslash \overline{A({\mathbb{A}})}^1}
\sum_{\gamma \in A({\mathbb{Q}})_\delta \backslash A({\mathbb{Q}})}
\overline{f}(x^{-1}\gamma ^{-1}\delta \gamma x)dx \\ 
=\sum_{\{\delta \}}
\int_{A({\mathbb{Q}})_\delta \backslash \overline{A({\mathbb{A}})}^1}
\overline{f}(x^{-1}\delta x)dx \\ 
=\sum_{\{\delta \}}
vol(A({\mathbb{Q}})_\delta \backslash \overline{A({\mathbb{A}})}_\delta ^1)
\sum_{x\in \overline{A({\mathbb{A}})}_\delta ^1
\backslash \overline{A({\mathbb{A}})}^1}
\overline{f}(x^{-1}\delta x),
\end{array}
\end{equation}
where $\{\delta \}$ runs over a complete set of representatives of
$\overline{A({\mathbb{A}})}^1$-conjugacy classes of $A({\mathbb{Q}})$.

This proves the

\begin{theorem}\label{thm:3.9}
(trace formula on $\overline{A({\mathbb{A}})}^1$):
If $\overline{f}\in {\mathcal{S}}(\overline{A({\mathbb{A}})}^1)$ then 
\[
\sum_{\overline{\chi }\in \Pi_u (\overline{A({\mathbb{A}})}^1)}
m(\overline{\chi })
tr(\overline{\chi }(\overline{f}))
=\sum_{\{\delta \}}
vol(A({\mathbb{Q}})_\delta \backslash \overline{A({\mathbb{A}})}_\delta ^1)
\sum_{x\in \overline{A({\mathbb{A}})}_\delta ^1
\backslash \overline{A({\mathbb{A}})^1}}
\overline{f}(x^{-1}\delta x),
\]
where $\{\delta \}$ runs over a complete set of representatives of
$\overline{A({\mathbb{A}})}^1$-conjugacy classes of
$A({\mathbb{Q}})$.
\end{theorem}

\subsection{proof of the correspondence}\label{ssec:rank0corr}

Let $\delta\in A({\mathbb{Q}})$.

We now prove the correspondence theorem mentioned at the beginning of this
subsection. To do this we need to relate the above twisted trace
formula on $\overline{A({\mathbb{A}})}^1$ to the ordinary trace
formula (Poisson summation) on $A({\mathbb{A}})^N$. This boils down to
relating the metaplectic orbital integrals 
\[
\int_{\overline{A({\mathbb{A}})}_\delta ^1
\backslash \overline{A({\mathbb{A}})^1}}
\overline{f}(x^{-1}\delta x)dx
\]
on $\overline{A({\mathbb{A}})}^1$ to ordinary orbital integrals
on $A({\mathbb{A}})^N$ (which are ``trivial'' since
$\overline{A({\mathbb{A}})}$ is abelian).

Let $\overline{f}$ be factorizable: 
\[
\overline{f}(x)=\prod_{p\leq \infty}\overline{f}_p(x_p),
\;\;\;\;\;\overline{f}_p\in {\mathcal{S}}(\overline{A({\mathbb{Q}}_p)}),
\]
so 
\begin{equation}\label{eq:3.15}
\int_{\overline{A({\mathbb{A}})}_\delta ^1
\backslash \overline{A({\mathbb{A}})^1}}
\overline{f}(x^{-1}\delta x)dx
=\prod_{p\leq \infty}
\sum_{x_p\in {\mathbb{Q}}_{p,\delta}^\times
\backslash {\mathbb{Q}}_p^\times}
\overline{f}_p(\delta ,(\delta,x_p)_n^2s_{\mathbb{A}}(\delta)^{-1} ),
\end{equation}
where
\[
\begin{array}{c}
{\mathbb{Q}}_{p,\delta}^\times
=\{x_p\in {\mathbb{Q}}_p^\times
\ |\ (x_p,\delta )_n=1\},\\
(x_p,\delta)_n=\prod_{v|p}(x_v,\delta_v)_{v,n},
\end{array}
\]
for $x_p=\prod_{v|p}x_v,\ \delta=\prod_{v|p}\delta_v$.
 
Let 
\begin{equation}\label{eq:3.16}
f_p(y)=
\sum_{x\in {\mathbb{Q}}_{p,y}^\times
\backslash {\mathbb{Q}}_p^\times}
\overline{f}_p(y ,(y,x)_n^2s_{\mathbb{A}}(y)^{-1} ),
\end{equation}
for $y\in \prod_{v|p}k_v^\times\cong A({\mathbb{Q}}_p)$.

\begin{remark}\label{rem:3.10}
Let $(h(t),1)\in \overline{A({\mathbb{Q}}_p}$ be the element
such that $h(t)$ is diagonal with diagonal entries $t,t^{-1}$.
For $y\in {\mathbb{Q}}_p^\times$, we have
$(h(y),1)^{(h(x),1)}=
(h(y) ,(y,x)_n^2) $.
Therefore, if $\overline{f}_p$ is genuine then
the sum in \eqref{eq:3.16} is zero unless
$y\in (\prod_{v|p}k_v^\times)^N\cong A({\mathbb{Q}}_p)^N$.
\end{remark}

\begin{lemma}\label{lem:3.11}
One may choose Haar measures so that

(a)\ If $p<\infty $\ then the map $\overline{f}_p\longmapsto f_p$
defines a map ${{\mathcal{S}}}(\overline{A({{\mathbb{Q}}}_p)})\rightarrow
{{\mathcal{S}}}(A({{\mathbb{Q}}}_p))$.

(b)\ If $p=\infty$ then the map
$\overline{f}_p\longmapsto f_p$\ defines a map
${{\mathcal{S}}}(\overline{A(\mathbb{R})})
\rightarrow {{\mathcal{S}}}(A({\mathbb{R}}))$.

(c) If $p<\infty $ and if $\overline{f}_p$
is a unit in the convolution
algebra ${\mathcal{S}}(\overline{A({\mathbb{Q}}_p)})$\ and
if $\overline{f}_p\longmapsto f_p$ then $f_p$
is the unit element of
the convolution algebra ${\mathcal{S}}(A({\mathbb{Q}}_p)^N)$.
\end{lemma}

For $\overline{f}\in {\mathcal{S}}(\overline{A({\mathbb{A}})^1})$
factorizable, let $\overline{f}\longmapsto f$ denote the map
obtained by taking the
product of the maps $\overline{f}_p\longmapsto f_p$.
By the above lemma, this defines a map 
\begin{equation}\label{eq:3.17}
{\mathcal{S}}(\overline{A({\mathbb{A}})}^1)\rightarrow {\mathcal{S}}(C_{{\mathbb{A}}}).
\end{equation}

As a corollary to the above lemma, we obtain

\begin{lemma}\label{lem:3.12}
Let\ $\overline{f}\in {\mathcal{S}}(\overline{A({\mathbb{A}})}^1)$ be a
factorizable function and define
$f\in {\mathcal{S}}(C_{\mathbb{A}})$ as above via
$\overline{f}\longmapsto f$. Then 
\[
\int_{\overline{A({\mathbb{A}})}_\delta ^1
\backslash \overline{A({\mathbb{A}})^1}}\overline{f}(x^{-1}\delta x)dx
=f(\delta ),
\]
if $\delta\in A({\mathbb{A}})^N$ and
\[
\int_{\overline{A({\mathbb{A}})}_\delta ^1
\backslash \overline{A({\mathbb{A}})^1}}\overline{f}(x^{-1}\delta x)dx
=0,
\]
if $\delta\in A({\mathbb{A}})-A({\mathbb{A}})^N$.
\end{lemma}

Note that if $\delta\in A({\mathbb{Q}})^N$ then
$\overline{A({\mathbb{A}})^1}_\delta=\overline{A({\mathbb{A}})^1}$.
We choose global measures so that
\[
vol(A({\mathbb{Q}}) \backslash \overline{A({\mathbb{A}})^1})
=vol(A({\mathbb{Q}})^N \backslash A({\mathbb{A}})^{1,N}).
\]
This and the above lemma imply that 
\begin{equation}\label{eq:3.18}
\begin{array}{c}
\sum_{\{\delta \}}
vol(A({\mathbb{Q}})_\delta \backslash \overline{A({\mathbb{A}})^1}_{\delta })
\int_{\overline{A({\mathbb{A}})^1}_{\delta }\backslash 
\overline{A({\mathbb{A}})^1}}
\overline{f}(x^{-1}\delta x)dx \\ 
=vol(A({\mathbb{Q}})^N \backslash A({\mathbb{A}})^{1,N})
\sum_{\gamma \in A({\mathbb{Q}})^N}
f(\gamma ),
\end{array}
\end{equation}
by equation \eqref{eq:3.16} and the remark above.

In this case, we have, by the Poisson summation formula for
$A({\mathbb{Q}})^N\backslash A({\mathbb{A}})^{1,N}$, we have 
\begin{equation}\label{eq:3.19}
\sum_{\overline{\chi }\in
\Pi_u (\overline{A({\mathbb{A}})^1})}m(\overline{\chi })
tr(\overline{\chi }(\overline{f}))
=\sum_{\chi \in \Pi (A({\mathbb{A}})^{1,N})}
f^{\wedge }(\chi ),
\end{equation}
for each factorizable
$\overline{f}\in {\mathcal{S}}(\overline{A({\mathbb{A}})}^1)$
which is genuine and $f\in {\mathcal{S}}(A({\mathbb{A}})^{1,N})$
is defined as above
via $\overline{f}\longmapsto f$.

Our correspondence theorem follows immediately from the two claims below:

\textsc{Claim 1}: For each fixed $\overline{\chi }$
$\in \Pi (\overline{A({\mathbb{A}})}^1)$ there is a sequence of
functions
$\overline{f}_m
=\overline{f}_{\overline{\chi },m}$ such that 
\[
\lim_{n\rightarrow \infty }
tr(\overline{\chi ^{\prime }}(\overline{f}_m))
=\left\{ 
\begin{array}{c}
1,\;\;\overline{\chi ^{\prime }}=\overline{\chi }, \\ 
0,\;\mathrm{otherwise}.
\end{array}
\right. 
\]

\textsc{Claim 2}: For each fixed
$\overline{\chi }\in \Pi (\overline{A({\mathbb{A}})}^1)$
there is a unique $\chi \in \Pi (A({\mathbb{A}})^{1,N})$ such that
$\overline{f}_m\longmapsto f_n=f_{\chi ,m}$ and 
\[
\lim_{m\rightarrow \infty }f_m^{\wedge }(\chi ^{\prime })=\left\{ 
\begin{array}{c}
1,\;\;\;\chi ^{\prime }=\chi , \\ 
0,\;\mathrm{otherwise}.
\end{array}
\right. 
\]

Claim 1 is an exercise in harmonic analysis on $A({\mathbb{A}})^1$ since
$tr(\overline{\chi ^{\prime }}(\overline{f}_m))
=\overline{f}_m^{\wedge }(\overline{\chi ^{\prime }})$.

To prove claim 2 we fix a genuine factorizable function $f$ on $A({\mathbb{A}})$
and work locally. If 
\[
\overline{\chi }=
Ind_{\overline{A({\mathbb{Q}}_p)^{1,N}}}^{\overline{A({\mathbb{Q}}_p)^1}}
(\chi )\in \Pi_u (\overline{A({\mathbb{A}})}^1)
\]
then we have 
\begin{equation}\label{eq:3.20}
\begin{array}{c}
\overline{f}^{\wedge }(\overline{\chi })
=\int_{\overline{A({\mathbb{Q}}_p)}}
\overline{f}(y)tr(\overline{\chi }(y))dy \\ 
=\sum_{h\in \overline{A({\mathbb{Q}}_p)^N}
\backslash \overline{A({\mathbb{Q}}_p)}}
\int_{\overline{A({\mathbb{Q}}_p)^N}}
\overline{f}(a)\chi ^h(a)da \\ 
=\sum_{h\in \overline{A({\mathbb{Q}}_p)^N}
\backslash \overline{A({\mathbb{Q}}_p)}}
\int_{\overline{A({\mathbb{Q}}_p)^N}}
\overline{f}(h^{-1}ah)\chi (a)da \\ 
=\int_{A({\mathbb{Q}}_p)^N\times 1}
f(a)\chi (a)da \\ 
=f^{\wedge }(\chi ).
\end{array}
\end{equation}

Claim 2 follows.

The 1-1 correspondence between $\Pi (\overline{A({\mathbb{A}})^1})$ and
$\Pi (A({\mathbb{A}})^{1,N})=\Pi (\overline{A({\mathbb{A}})^{1,N}})_{gen}$
is given as follows. The map 
\[
\Pi (\overline{A({\mathbb{A}})^{1,N}})
\rightarrow \Pi (\overline{A({\mathbb{A}})}^1)
\]
is the induction map. In fact, the map
$\Pi (\overline{A({\mathbb{A}})^1})\rightarrow
\Pi (A({\mathbb{A}})^{1,N})$
is given locally by the map 
\[
\overline{\chi }
=Ind_{\overline{A({\mathbb{Q}}_p)^N}}^{\overline{A({\mathbb{Q}}_p)}}
(\chi )\longmapsto \chi 
\]
This proves the desired correspondence. $\Box $

Let ${\mathcal{O}}$ denote a (singleton) conjugacy
class in $A({\mathbb{Q}})$ and let
$\chi\in \Pi_u(\overline{A({\mathbb{A}})})$ denote a genuine irreducible
representation. Define $I^A_{\mathcal{O}}$ by
\begin{equation}\label{eq:3.21}
I^A_{\mathcal{O}}(f)=
\sum_{\delta\in {\mathcal{O}}}
f(\delta)
\end{equation}
and $I^A_{\chi}$ by
\begin{equation}\label{eq:3.22}
I^A_{\chi}(f)=
{\rm tr}(\chi(f)) 
\end{equation}
for $f\in {\mathcal{H}}(\overline{A({\mathbb{A}})})$, where
${\mathcal{H}}(\overline{A({\mathbb{A}})})$ is as in \S\ref{sssec:defs} below.


